% =====================================================================
%  Etude theorique complete - niveau prepa, schemas clairs avec axes.
%  Modele : poutre encastree (bord de l'ecran). Compilable seul.
% =====================================================================
\documentclass[11pt,aspectratio=169]{beamer}
\usetheme{Madrid}\usecolortheme{whale}
\setbeamertemplate{navigation symbols}{}
\usepackage[utf8]{inputenc}
\usepackage[T1]{fontenc}
\usepackage[french]{babel}
\usepackage{amsmath,amssymb}
\usepackage{tikz}
\usetikzlibrary{arrows.meta,calc}
\definecolor{verre}{RGB}{30,120,180}
\definecolor{danger}{RGB}{200,40,40}
\definecolor{ok}{RGB}{30,150,90}
\definecolor{gristik}{RGB}{120,120,120}

\begin{document}
\section{Étude théorique}

% ---------- S1 : elasticite ----------
\begin{frame}{1. Élasticité linéaire : les grandeurs de base}
  \begin{columns}[c]
  \column{0.40\textwidth}
    \centering
    \begin{tikzpicture}[scale=1]
      \draw[verre,thick,fill=verre!12](0,0) rectangle (2.6,0.7);
      \draw[-{Latex},danger,thick](0,0.35)--(-0.9,0.35) node[left,font=\small]{$F$};
      \draw[-{Latex},danger,thick](2.6,0.35)--(3.5,0.35) node[right,font=\small]{$F$};
      \draw[{Latex}-{Latex}](0,-0.35)--(2.6,-0.35) node[midway,below,font=\footnotesize]{$L$};
      \node[font=\tiny] at (1.3,1.05){section $S$};
    \end{tikzpicture}
  \column{0.60\textwidth}
    \footnotesize
    \begin{itemize}\setlength{\itemsep}{7pt}
      \item \textbf{Contrainte} $\sigma=\dfrac{F}{S}$ : force par unité de
            surface (Pa).
      \item \textbf{Déformation} $\varepsilon=\dfrac{\Delta L}{L}$ : allongement
            relatif.
      \item \textbf{Loi de Hooke} : $\sigma=E\,\varepsilon$ ($E$ : module
            d'Young).
      \item \textbf{Poisson} : $\varepsilon_{\perp}=-\nu\,\varepsilon$.
    \end{itemize}
    \scriptsize Verre : $E\approx 70$~GPa, $\nu\approx 0{,}22$.
  \end{columns}
\end{frame}

% ---------- S2 : modele (extraction) ----------
\begin{frame}{2. Le modèle : on isole une poutre de l'écran}
  \centering
  \begin{tikzpicture}[scale=0.85]
    % ecran
    \draw[verre,thick,fill=verre!8,rounded corners=3pt](0,0) rectangle (1.7,3.2);
    \fill[danger!30](0.65,0) rectangle (1.0,3.2);
    \node[font=\scriptsize] at (0.85,3.5){écran};
    % fleche
    \draw[-{Latex},thick](2.0,1.6)--(2.9,1.6);
    \node[font=\scriptsize] at (2.45,1.95){on isole};
    % poutre longue et fine
    \begin{scope}[shift={(3.3,1.5)}]
      \fill[verre!18](0,0) rectangle (5,0.2);
      \draw[verre,thick](0,0) rectangle (5,0.2);
      \fill[verre!28](0,0.2)--(5,0.2)--(5.3,0.42)--(0.3,0.42)--cycle;
      \draw[verre,thick](0,0.2)--(5,0.2)--(5.3,0.42)--(0.3,0.42)--cycle;
      \draw[verre,thick](5,0)--(5.3,0.22)--(5.3,0.42);
      \draw[verre,thick](5,0.2)--(5.3,0.42);
      \draw[{Latex}-{Latex}](0,-0.3)--(5,-0.3) node[midway,below,font=\small]{$L$};
      \draw[{Latex}-{Latex}](-0.28,0)--(-0.28,0.2) node[midway,left,font=\small]{$e$};
      \draw[{Latex}-{Latex}](5.12,-0.05)--(5.45,0.18) node[midway,right,font=\small]{$b$};
    \end{scope}
  \end{tikzpicture}
  \vspace{0.4cm}
  \par\footnotesize On isole une \textbf{fine poutre} ($e\ll L$) :
  $L\sim$ quelques cm, $b\sim 1$~cm, $e\approx 0{,}5$~mm.
\end{frame}

% ---------- S3 : flexion et ligne neutre ----------
\begin{frame}{3. Sous une charge : flexion et ligne neutre}
  \begin{columns}[c]
  \column{0.58\textwidth}
    \centering
    \begin{tikzpicture}[scale=1]
      % axes
      \draw[-{Latex}](-0.3,0)--(6.0,0) node[right,font=\small]{$x$};
      \draw[-{Latex}](0,-1.2)--(0,1.1) node[above,font=\small]{$y$};
      % bande flechie (epaisse pour la lisibilite)
      \fill[verre!15] (0,0.25) .. controls (2.75,-0.35) .. (5.5,0.25)
        -- (5.5,-0.25) .. controls (2.75,-0.85) .. (0,-0.25) -- cycle;
      \draw[verre,thick](0,0.25) .. controls (2.75,-0.35) .. (5.5,0.25);
      \draw[verre,thick](0,-0.25) .. controls (2.75,-0.85) .. (5.5,-0.25);
      % ligne neutre
      \draw[danger,dashed,thick](0,0) .. controls (2.75,-0.60) .. (5.5,0);
      % charge
      \draw[-{Latex},danger,thick](2.75,1.0)--(2.75,-0.35) node[pos=0.15,right,font=\tiny]{$F$};
      % legendes (hors du trace)
      \node[ok,font=\tiny] at (1.2,0.7){fibres comprimées};
      \node[danger,font=\tiny] at (4.3,-1.05){fibres tendues};
      \node[danger,font=\tiny,right] at (5.55,-0.05){ligne neutre};
    \end{tikzpicture}
  \column{0.42\textwidth}
    \footnotesize
    La poutre fléchit :
    \begin{itemize}\setlength{\itemsep}{5pt}
      \item un côté se \textbf{comprime}, l'autre se \textbf{tend} ;
      \item la \textbf{ligne neutre} (au milieu) garde sa longueur.
    \end{itemize}
    Une fibre à la distance $y$ subit :
    \[ \sigma_{xx}=-E\,\frac{y}{R}. \]
  \end{columns}
\end{frame}

% ---------- S4 : efforts internes ----------
\begin{frame}{4. Efforts internes : effort tranchant et moment}
  \begin{columns}[c]
  \column{0.46\textwidth}
    \centering
    \begin{tikzpicture}[scale=1]
      \draw[-{Latex}](-0.3,0)--(5.0,0) node[right,font=\small]{$x$};
      \draw[-{Latex}](0,-1.0)--(0,1.0) node[above,font=\small]{$y$};
      \draw[verre,very thick](0,0)--(4.3,0);
      \draw[gristik,dashed](2.4,0.9)--(2.4,-0.9);
      \node[font=\tiny,below] at (2.4,-0.95){section en $x$};
      \draw[-{Latex},danger,thick](2.4,-0.1)--(2.4,-0.8) node[right,font=\tiny]{$V(x)$};
      \draw[-{Latex},ok,thick](2.1,0.34) arc (160:20:0.3);
      \node[ok,font=\tiny] at (2.4,0.78){$M(x)$};
    \end{tikzpicture}
  \column{0.54\textwidth}
    \footnotesize
    On coupe la poutre en $x$ et on écrit l'équilibre de la portion :
    \[ \textstyle\sum\vec F=\vec 0,\qquad \sum\vec{\mathcal M}=\vec 0. \]
    On en déduit :
    \begin{itemize}\setlength{\itemsep}{4pt}
      \item l'\textbf{effort tranchant} $V(x)$ (force, en N) ;
      \item le \textbf{moment fléchissant} $M(x)$ (en N$\cdot$m).
    \end{itemize}
  \end{columns}
\end{frame}

% ---------- S5 : M(x) ----------
\begin{frame}{5. Calcul du moment fléchissant $M(x)$}
  \begin{columns}[c]
  \column{0.5\textwidth}
    \centering
    \begin{tikzpicture}[scale=0.9]
      % poutre encastree + F
      \draw[thick](0,1.7)--(0,2.9);
      \foreach \yy in {1.8,2.05,...,2.8}{\draw[gristik](0,\yy)--(-0.18,\yy+0.14);}
      \draw[verre,very thick](0,2.3)--(4,2.3);
      \draw[-{Latex},danger,thick](4,3.0)--(4,2.35) node[midway,right,font=\tiny]{$F$};
      \node[font=\tiny,below] at (4,2.2){$x=L$};
      % diagramme M(x)
      \draw[-{Latex}](0,0)--(4.7,0) node[right,font=\tiny]{$x$};
      \draw[-{Latex}](0,0)--(0,1.5) node[above,font=\tiny]{$M$};
      \draw[verre,very thick](0,1.2)--(4,0);
      \node[font=\tiny,left] at (0,1.2){$FL$};
      \node[font=\tiny,below] at (4,0){$L$};
    \end{tikzpicture}
  \column{0.5\textwidth}
    \footnotesize
    Équilibre global : $R_y=-F$, $\;M_e=-LF$.
    \medskip\\
    Portion $[0,x]$ :
    \[ V_y(x)=F,\quad M_e+xV_y(x)+M(x)=0 \]
    \[ \boxed{\;M(x)=F\,(L-x)\;} \]
    Maximal à l'encastrement : $M_{\max}=FL$.
  \end{columns}
\end{frame}

% ---------- S6 : contrainte locale ----------
\begin{frame}{6. Contrainte locale dans une section}
  \begin{columns}[c]
  \column{0.5\textwidth}
    \centering
    \begin{tikzpicture}[scale=1]
      % section : axe y vertical
      \draw[verre,thick,fill=verre!10](0,-1.4) rectangle (0.55,1.4);
      \draw[-{Latex}](0,-1.7)--(0,1.7) node[above,font=\small]{$y$};
      \draw[gristik,dashed](-0.4,0)--(2.6,0) node[right,font=\tiny]{ligne neutre};
      % profil lineaire (fleches)
      \foreach \z/\l in {0.5/0.5,1.0/1.0,1.4/1.4}{
        \draw[-{Latex},ok](0.55+\l,\z)--(0.55,\z);}
      \foreach \z/\l in {0.5/0.5,1.0/1.0,1.4/1.4}{
        \draw[-{Latex},danger](0.55,-\z)--(0.55+\l,-\z);}
      \node[ok,font=\tiny,right] at (1.6,1.4){comprimé};
      \node[danger,font=\tiny,right] at (1.6,-1.4){tendu};
    \end{tikzpicture}
  \column{0.5\textwidth}
    \footnotesize
    La contrainte varie \textbf{linéairement} avec $y$ :
    \[ \sigma_{xx}(y)=-E\,\frac{y}{R}. \]
    Elle est \textbf{nulle} sur la ligne neutre et \textbf{maximale en surface}
    ($y=\pm e/2$) :
    \[ \sigma_{\max}=E\,\frac{e}{2}\,\frac{1}{R}. \]
  \end{columns}
\end{frame}

% ---------- S7 : M = EI/R ----------
\begin{frame}{7. Du moment à la courbure : $M=EI/R$}
  \begin{columns}[c]
  \column{0.42\textwidth}
    \centering
    \begin{tikzpicture}[scale=1]
      \draw[verre,thick,fill=verre!10](0,-1.4) rectangle (0.55,1.4);
      \draw[-{Latex}](0,-1.7)--(0,1.7) node[above,font=\small]{$y$};
      \draw[gristik,dashed](-0.3,0)--(2.2,0);
      \fill[danger!40](0,0.8) rectangle (0.55,0.95);
      \draw[{Latex}-{Latex}](0.7,0)--(0.7,0.87) node[midway,right,font=\tiny]{$y$};
      \draw[-{Latex},danger](1.0,0.87)--(1.8,0.87) node[right,font=\tiny]{$\mathrm{d}F$};
    \end{tikzpicture}
  \column{0.58\textwidth}
    \footnotesize
    Chaque fibre porte $\mathrm{d}F=\sigma_{xx}\,\mathrm{d}S$. En sommant les
    moments sur la section :
    \[ M=\int_{S}\frac{E\,y^{2}}{R}\,\mathrm{d}S=\frac{E\,I}{R}, \]
    avec le \textbf{moment quadratique}
    \[ I=\int_{S} y^{2}\,\mathrm{d}S. \]
  \end{columns}
\end{frame}

% ---------- S8 : equation de la deformee ----------
\begin{frame}{8. L'équation de la déformée $u_y(x)$}
  \begin{columns}[c]
  \column{0.54\textwidth}
    \centering
    \begin{tikzpicture}[scale=1]
      \draw[thick](0,-0.6)--(0,0.7);
      \foreach \yy in {-0.5,-0.25,...,0.6}{\draw[gristik](0,\yy)--(-0.2,\yy+0.16);}
      \draw[-{Latex}](0,0)--(5.5,0) node[right,font=\small]{$x$};
      \draw[-{Latex}](0,-1.9)--(0,1.0) node[above,font=\small]{$y$};
      \draw[gristik,dashed](0,0)--(5,0);
      \node[gristik,font=\tiny] at (4.2,0.22){repos};
      \draw[verre,very thick](0,0) .. controls (2.5,-0.8) and (3.5,-1.5) .. (5,-1.7);
      \node[verre,font=\tiny] at (4.0,-1.78){déformée};
      \draw[densely dotted](3.0,0)--(3.0,-1.02);
      \draw[{Latex}-{Latex}](3.25,0)--(3.25,-1.02) node[midway,right,font=\tiny]{$u_y(x)$};
      \draw (3.0,0.08)--(3.0,-0.08);
      \node[font=\tiny,above] at (3.0,0.10){$x$};
    \end{tikzpicture}
  \column{0.46\textwidth}
    \footnotesize
    $u_y(x)$ : \textbf{déplacement vertical} de la fibre neutre (la déformée).
    \medskip\\
    Courbure : $\dfrac{1}{R}\approx\dfrac{\mathrm{d}^{2}u_y}{\mathrm{d}x^{2}}$.
    Avec $M=EI/R$ :
    \[ \boxed{\;EI\,\frac{\mathrm{d}^{2}u_y}{\mathrm{d}x^{2}}=M(x)\;} \]
  \end{columns}
\end{frame}

% ---------- S9 : resolution ----------
\begin{frame}{9. Résolution : la déformée $u_y(x)$}
  \footnotesize
  On remplace $M(x)=F(L-x)$ :
  \[ EI\,\frac{\mathrm{d}^{2}u_y}{\mathrm{d}x^{2}}=F(L-x). \]
  Conditions d'encastrement : $u_y(0)=0$ et $u_y'(0)=0$. Deux intégrations donnent :
  \[ \boxed{\;u_y(x)=\frac{F}{EI}\left(\frac{Lx^{2}}{2}-\frac{x^{3}}{6}\right)\;} \]
  d'où la fléche maximale au bord libre :
  \[ u_{\max}=u_y(L)=\frac{FL^{3}}{3EI}. \]
\end{frame}

% ---------- S10 : sigma(u_max) + energie + graphe ----------
\begin{frame}{10. De la fléche à $\sigma_{\max}$, puis à $1/\sqrt{e}$}
  \begin{columns}[c]
  \column{0.5\textwidth}
    \footnotesize
    \textbf{Lien $\sigma_{\max}$–fléche} : avec $\sigma_{\max}=E\frac{e}{2}\frac{M_{\max}}{EI}$
    et $M_{\max}=\frac{3EI\,u_{\max}}{L^{2}}$ :
    \[ \sigma_{\max}=\frac{3Ee}{2L^{2}}\,u_{\max}. \]
    \textbf{Énergie} : $\tfrac12 k u_{\max}^{2}=\eta E_c$, $k=\dfrac{3EI}{L^{3}}\propto e^{3}$,
    donc $u_{\max}\propto e^{-3/2}$ et :
    \[ \boxed{\;\sigma_{\max}\propto 1/\sqrt{e}\;} \]
  \column{0.5\textwidth}
    \centering {\scriptsize $\sigma_{\max}$ vs épaisseur (chute de 1~m)}
    \begin{tikzpicture}[scale=0.95]
      \draw[-{Latex}](0,0)--(7.0,0) node[right,font=\scriptsize]{$e$ (mm)};
      \draw[-{Latex}](0,0)--(0,4.7) node[above,font=\scriptsize]{$\sigma_{\max}$ (MPa)};
      \foreach \e in {0.3,0.5,1.0,1.5}{
        \draw ({\e*4.2},0)--({\e*4.2},-0.1) node[below,font=\tiny]{\e};}
      \foreach \s/\yc in {300/1.5,600/3.0,900/4.5}{
        \draw (0,\yc)--(-0.1,\yc) node[left,font=\tiny]{\s};}
      \draw[danger,dashed](0,4.25)--(6.6,4.25) node[right,font=\tiny,danger]{850};
      \draw[verre,very thick,domain=0.28:1.55,samples=70,smooth]
        plot ({\x*4.2},{622*sqrt(0.5/\x)/200});
      \foreach \e/\s in {0.3/803,0.5/622,1.0/440,1.5/359}{
        \fill[danger] ({\e*4.2},{\s/200}) circle (1.4pt);
        \node[font=\tiny,above right] at ({\e*4.2},{\s/200}){\s};}
    \end{tikzpicture}
  \end{columns}
\end{frame}

\end{document}
